% Corrigé intégré automatiquement pour la banque Exercices.
% Source : A_Integrer/DDS_04/073_SLCI_Retard/073_SLCI_Retard.tex
\begin{corrigebox}[Corrigé]
\CorrigeQuestion{1}
$H'_{\text{frein}} ( p) = \dfrac{1,9}{1+3,3 p} $.
\CorrigeQuestion{2}
$H_{\text{frein}} ( p) = \dfrac{1,9}{1+3,3 p} e^{-1,7 p}$.
\CorrigeQuestion{3}
$M\varphi = 100\degres$ et $\omega_c = \SI{2}{rad.s^{-1}}$.
\CorrigeQuestion{4}
On factorise la fonction de transfert sous forme canonique. Chaque gain, intégrateur, zéro et pôle apporte sa pente et sa phase ; les contributions sont ensuite sommées. La marge de phase se lit à la pulsation où le gain vaut $0$ dB et la marge de gain à la pulsation où la phase vaut $-180^\circ$.
\CorrigeQuestion{5}
$M\varphi = -180\degres$ et $\omega_c = \SI{2}{rad.s^{-1}}$ et $MG = \SI{-4}{dB}$.
\CorrigeQuestion{6}
Le correcteur est introduit dans la boucle ouverte, puis ses paramètres sont choisis pour placer la pulsation de coupure et obtenir la marge de phase demandée. Pour un PI, le zéro $1/T_i$ est placé sous la pulsation de coupure ; le gain est ensuite réglé pour imposer $|L(j\omega_c)|=1$.
\end{corrigebox}
